\chapter{Oxygen, Sulfur, Nitrogen, and Phosphorus: Life and Environment's Framework}

\begin{kaichang}
Find matches, or look at the kitchen gas stove if none are available, and do something familiar but rarely considered: strike a match.

A rasp, a puff of white smoke, then flame and a pungent smell. Nearly every actor in that second belongs to this chapter. The box's rough red-brown strip contains \textbf{phosphorus}; the head contains \textbf{sulfur} and a white powder, potassium chlorate, that releases \textbf{oxygen} when heated. Friction ignites phosphorus, which ignites sulfur, which lights the wooden stem. Oxygen supplied by the head intensifies burning.

Now two questions. If fire burns in air, why does the head carry its own oxygen package? And air is mostly nitrogen, nearly four fifths. The whole match is immersed in it, so why does nitrogen stand aside, neither helping nor interfering?

Why do these four elements jointly support your bones and energy while also stirring up acid rain and lakes? Guess before reading on.
\end{kaichang}

\section{Observations: One Encourages Fire, One Watches}

Begin with observable phenomena alone.

You may have seen a glowing splint inserted into an oxygen jar in class, suddenly reigniting brightly. A candle burns more vigorously in pure oxygen than air; thin iron wire that will not burn in air throws sparks in oxygen. Conversely, cover a candle with a sealed glass bell and it weakens and goes out as the combustion-supporting component is used.

Nitrogen shows the opposite. Puffy chip bags contain nitrogen, displacing oxygen to prevent rancidity and cushioning against breakage. Liquid nitrogen presents another face: about $-196\,^{\circ}\mathrm{C}$, cold enough to make petals shatter when tapped, yet useful for preserving cells and seeds because it remains nearly unreactive despite the cold.

Rust is oxygen in slow motion. A nail in a damp corner develops red rust over weeks, a nearly imperceptible reaction consuming something in the air. The following experiment makes that consumption visible.

\begin{shiyan}
\textbf{Watching a Portion of Air Disappear}

\textbf{Materials}: a small wad of hardware-store steel wool, a narrow-necked glass bottle, a shallow dish, water, and a few drops of white vinegar.

\textbf{Procedure}: Wet the wool with vinegar to remove protective surface oil and push it into the bottle's bottom. Fill the dish with water, invert the bottle into it, and equalize water levels inside and outside. Mark the level. Leave by a window and observe daily for three to five days.

\textbf{What to observe}: The wool turns reddish brown and the bottle's water rises slowly, eventually occupying roughly one fifth of the original air volume.

\textbf{Safety}: No fire is needed. Do not try any method of making pure oxygen to accelerate it. Steel wool has sharp edges; wash hands afterward and keep vinegar out of eyes.

Those centimeters of rising water fill space left by oxygen moving from air into rust. At home, you have roughly measured air's oxygen share, about one fifth.
\end{shiyan}

\section{Particles: Oxygen Holds Two Hands, Nitrogen Three}

Why are air's two major components active and idle respectively? Zoom into molecules.

Chapter 10's neighborhoods place oxygen and sulfur in Group 16 with 6 outer electrons, nitrogen and phosphorus in Group 15 with 5. Bringing in one or two electrons to fill the shell more often lowers energy. This is an energy rule, not a wish. All four therefore have strong nonmetallic tendencies, sharing or taking electrons, a prerequisite for forming life's framework.

Atmospheric oxygen is \chem{O_2}, two atoms sharing two electron pairs like two firmly clasped hands (Chapter 19 explains covalent bonds). Nitrogen is also diatomic, \chem{N_2}, but shares \textbf{three pairs}. That extra hand matters: separating nitrogen atoms costs about \ud{945}{kJ\,mol^{-1}}, nearly twice oxygen's \ud{498}{kJ\,mol^{-1}}, among the strongest common diatomic bonds.

The macroscopic contrast now has a particle explanation. Nitrogen's startup fee is too expensive, rather than its being unwilling. Typical room-temperature collisions cannot pay, blocking reaction. This concerns \textbf{speed}, not direction; Chapters 27 and 29 separate possibility from rate. Oxygen's fee is much lower, affordable to wood, iron wire, and bodily glucose.

Now the symbols. Simplifying wood's main component as carbon gives combustion:
\[\chem{C} + \chem{O_2} \rightarrow \chem{CO_2}\]
Rusting is slower but similar: iron surrenders electrons to oxygen and forms oxides. Your continuous respiration likewise records organic matter $+$ oxygen $\rightarrow$ carbon dioxide $+$ water, the same direction as burning, divided by life into dozens of gradual steps (Chapter 54).

\section{Oxygen Is Not Fuel: It Accepts Electrons}

Pause to dismantle a widespread claim.

\begin{wuqu}
``Oxygen burns'' or ``oxygen is fuel'' is wrong and genuinely hazardous. Fuel is the oxidized partner, such as wood, gasoline, or hydrogen. Oxygen \textbf{accepts electrons}, making it the oxidizing agent. Oxygen itself never burns: pure oxygen in a balloon cannot simply be ignited without fuel. The reverse is dangerous: pure oxygen makes ordinarily mild materials such as grease, clothing, and hair burn fiercely, hence bans on oil contamination and flames near hospital oxygen cylinders. Defining oxidation and reduction as oxygen gain and loss is only their historical starting point. Their essence is electron transfer, revisited in Chapter 33.
\end{wuqu}

Combustion thus needs fuel, oxygen, and sufficiently high starting temperature; removing any extinguishes it. Covering a candle cuts oxygen; water mainly cools. This half-answers the first opening question: heated potassium chlorate releases \chem{O_2}, creating oxygen \textbf{far more concentrated than air} at the tiny match head to ensure ignition. Air's 21\% often cannot act quickly enough on a tiny flame.

\begin{qiaoliang}
\textbf{How Much Oxygen Is in a Breath of Air?} Dry air is about 78\% nitrogen, 21\% oxygen, and 1\% mainly argon and carbon dioxide by volume. A classroom 9 by 7 by 3 meters has volume $189\ \mathrm{m^3}$ and about $189\times 21\% \approx 40\ \mathrm{m^3}$ oxygen. A seated student consumes roughly $0.25\ \mathrm{L}$ per minute, $0.36\ \mathrm{m^3}$ daily, so theoretically one person could use the classroom's oxygen for over a hundred days. Why does it feel stuffy? Often \textbf{rising carbon dioxide}, not oxygen shortage, causes discomfort. Exhalation raises carbon dioxide from about 0.04\% to 4\%, and its accumulation in blood is the main stuffiness signal. Feeling and measurement differ, which is why instruments beat intuition.
\end{qiaoliang}

\section{Ozone: One Molecule, Two Reputations}

Oxygen also has a three-atom form, ozone \chem{O_3}. Less structurally stable than \chem{O_2}, it is more reactive and a stronger oxidizer. Its reputation depends entirely on location.

At 20--30 kilometers above Earth, ozone forms a thin layer. Compressed to ordinary ground-level temperature and pressure, it would be only about 3 millimeters thick. Yet it absorbs most of sunlight's most biologically damaging ultraviolet. Without it, sunlight would break DNA bonds directly and terrestrial life would struggle to exist. At this height, ozone is Earth's sunshade.

At breathing level, ozone is pollution. Strong summer sunlight drives reactions among vehicle and industrial nitrogen oxides and volatile organic compounds, forming ozone and other irritants as photochemical smog. Ground-level ozone irritates eyes and airways, especially endangering asthma patients, and harms leaves. One molecule, protective above and irritating below: \textbf{location determines role}.

This is a key to environmental chemistry. Before calling a substance good or bad, specify where and at what concentration. Disinfection cabinets and printers also produce ozone with a faint thunderstorm smell. Lightning really does turn some \chem{O_2} into \chem{O_3}, but that concentration is far too low to purify indoor air, and concentrated ozone from generators harms people too.

In the later twentieth century, humanity found an unexpected effect: widely used chlorofluorocarbon refrigerants reaching the upper atmosphere release chlorine atoms under ultraviolet. One chlorine atom repeatedly catalyzes destruction of tens of thousands of ozone molecules, opening the Antarctic ozone hole. Later discussions of catalysts and radicals unpack this history. Here remember that an apparently inert molecule, moved upward, can become scissors cutting the sunshade.

\section{Nitrogen: A Tough Triple Bond and an Industry Feeding Half the World}

The second opening answer is now clear: \chem{N_2}'s triple bond is so strong that a match's temperature cannot pay its splitting fee.

Hard to break does not mean impossible. Lightning's momentary heat breaks the bond and combines nitrogen with oxygen into oxides carried by rain into soil, natural nitrogen fertilizer. Bacteria in legume root nodules use enzymes at room temperature to turn \chem{N_2} stepwise into ammonia, biological nitrogen fixation. Around Chapter 54 we revisit these microscopic factories. Natural fixation, however, cannot keep pace with farmland's nitrogen demand under population growth.

Why is nitrogen vital? Proteins and DNA require it, but plants absorb ready-made nitrogen ions, nitrates and ammonium, rather than atmospheric \chem{N_2}. They are people beside a granary without its key. Tons of nitrogen float overhead while crops lack it below. Whoever fixes atmospheric nitrogen into usable form controls a food-production valve.

\begin{zhengju}
Late nineteenth-century European farming increasingly depended on Chilean nitrate shipped by sea. In a famous 1898 speech, William Crookes warned that consumption would exhaust deposits within decades, wheat would fail to support population, and bread-eating civilization would starve. The problem was plain: convert inexhaustible atmospheric \chem{N_2} to ammonia.

Chemistry already revealed conflicting accounts. Ammonia synthesis releases heat, so lower temperature favors ammonia at equilibrium. But low temperature makes triple-bond breaking uselessly slow. Nernst therefore considered the industrial route impractical. Haber disagreed and made the crucial move: \textbf{take the reaction to high pressure}. Compression crowds gas molecules and shifts equilibrium toward fewer product molecules, forcing yield upward. In 1909, around 200 atmospheres and $500\,^{\circ}\mathrm{C}$ with an iron-containing catalyst, his apparatus continuously converted nitrogen and hydrogen into ammonia. Yield data overturned Nernst's pessimism; Nernst verified them and publicly admitted error.

An apparatus is not a factory. Pressure-resistant steel vessels, inexpensive hydrogen, and long-lived catalysts were separate engineering challenges. Bosch's BASF team spent over four years and tested more than twenty thousand catalyst formulations before opening the first ammonia plant in 1913.

The ending is not simple. Ammonia synthesis freed fertilizer from mineral limits and now supplies nitrogen feeding about half humanity. The same equipment, however, produced nitric acid for explosives in both world wars. Haber was eventually expelled from the country he served because he was Jewish. A technology entering society never brings only its original promise; Volume III repeatedly returns to this theme.
\end{zhengju}

The ammonia account occupies one line:
\[\chem{N_2} + \chem{3H_2} \rightleftharpoons \chem{2NH_3}\]
The reversible arrow signals a two-way equilibrium (Chapter 31). Industry continuously removes ammonia and compresses unreacted gases back through the process, turning an idle gas into production-line feedstock.

\begin{qiaoliang}
\textbf{How Much Haber's Nitrogen Is in You?} Global ammonia synthesis fixes about 120 million tonnes of nitrogen annually. Divided among 8 billion people, that is 15 kilograms each, equivalent to the nitrogen in about 90 kilograms of ammonium-nitrate fertilizer yearly. A 60-kilogram person is about 3\% nitrogen, 1.8 kilograms, almost entirely in proteins and DNA. Researchers estimate that roughly half of modern humanity's nitrogen atoms trace back to synthetic ammonia. About 0.9 kilograms of your nitrogen once emerged from \chem{N_2} in a high-pressure steel vessel. Elements follow no road signs; they flow through air, fertilizer, crops, and you.
\end{qiaoliang}

\section{Phosphorus: Stone in Bones and Energy Change in Cells}

Nitrogen's downstairs neighbor, phosphorus, also has 5 outer electrons in Group 15, but is solid and far less inclined to stay home.

Elemental phosphorus has two famous faces. White phosphorus is \chem{P_4}, four atoms forming a small tetrahedron. Tight angles strain bonds and raise energy, making it so reactive that it can ignite spontaneously in air around $40\,^{\circ}\mathrm{C}$. It must be stored underwater in laboratories; never handle it at home. Red phosphorus forms a larger, less-strained network and is milder, igniting only around $240\,^{\circ}\mathrm{C}$. Thus it safely coats matchboxes until friction supplies brief heat. The same element's different atomic connections produce utterly different temperaments. Allotropy will shine again when Chapter 36 discusses carbon.

Life uses phosphorus more calmly. Over eighty percent of bones' and teeth's inorganic content is calcium hydroxyphosphate, containing calcium, phosphate, and hydroxide. Bone is essentially stone framework reinforced with collagen. Cell-membrane phospholipids contain phosphate; DNA's helical handrails alternate sugars and phosphates (Chapter 58). Phosphorus's most famous role hides in three letters: ATP.

ATP, adenosine triphosphate, is often called cellular energy currency. The familiar claim that energy is stored in high-energy phosphate bonds and released by breaking them is \textbf{wrong}. Remember: \textbf{breaking any chemical bond absorbs energy; forming new bonds releases it} (Chapter 27). In ATP hydrolysis,
\[\chem{ATP} + \chem{H_2O} \rightarrow \chem{ADP} + \chem{Pi}\]
the products ADP and phosphate are much more stable than the reactants. Crowded negative ATP charges separate, reducing repulsion, and products disperse through water in more ways. Old-bond breaking costs energy, but new bonds and charge dispersal return more. The \textbf{net profit} supplies usable energy. ATP resembles spendable small change, not a packet of gunpowder: hydrolysis is profitable overall.

Environmental phosphorus has another face: \textbf{eutrophication}. Relatively scarce in the crust and scarcely entering air, phosphorus is often algae's most limiting food in lakes and rivers. Other nutrients cannot trigger a bloom without it; add phosphorus and growth explodes. Phosphate detergents and fertilizer runoff once turned many waters paint-green. Blooming algae die, bacterial decomposition consumes oxygen, and fish and other animals suffocate. Hence many countries banned phosphate laundry detergents. Nitrogen and phosphorus cycling through rock, soil, organisms, and water forms a vital ecological ledger, fully unfolded in Chapter 16.

\section{Sulfur: Rotten Eggs, Acid Plants, and Sour Rain}

Oxygen's downstairs neighbor sulfur also has 6 outer electrons but one extra shell, subtly changing its character. It takes electrons less strongly and can therefore perform jobs oxygen cannot.

Elemental sulfur is a brittle yellow solid, most stable at room temperature as eight-atom crown rings, \chem{S_8}. Its relationship with life begins in smells: rotten eggs, spoiled cabbage, and pungent volcanic odors warn of small sulfur molecules. Rotten eggs' main actor is \chem{H_2S}, produced as bacteria decompose sulfur-containing amino acids. Hair and keratin contain many disulfide bridges; perm chemicals break and reconnect these. Remember that concentrated \chem{H_2S} is poisonous and disables smell, so absence of odor does not establish safety. Warnings near sewers and biogas pits must be taken seriously; professionals alone should handle these situations.

Sulfur's industrial importance is captured by sulfuric-acid production as a nation's industrial thermometer. Fertilizers, batteries, metallurgy, synthetic fibers, and medicines all use it. Modern production burns sulfur or sulfur ore to \chem{SO_2}, oxidizes that on catalysts, commonly vanadium pentoxide, to \chem{SO_3}, then absorbs it in water:
\[\chem{S} + \chem{O_2} \rightarrow \chem{SO_2}\]
\[\chem{2SO_2} + \chem{O_2} \rightleftharpoons \chem{2SO_3}\]
\[\chem{SO_3} + \chem{H_2O} \rightarrow \chem{H_2SO_4}\]
The middle step is difficult because of its startup fee. Catalysts provide a cheaper path, changing rate but neither direction nor ultimate accounts (Chapter 30).

Direct chimney release of \chem{SO_2}, however, causes trouble. Coal and petroleum naturally contain sulfur, burned into sulfur dioxide that rises in flue gas. Air and cloud water oxidize and dissolve it gradually into dilute sulfuric acid falling back to Earth. Vehicle nitrogen oxides follow similar routes to nitric acid: acid rain. \textbf{Normal rain is already mildly acidic}, around pH 5.6, because dissolved atmospheric \chem{CO_2} makes some carbonic acid. Acid rain falls substantially below that background. It dissolves limestone sculptures and buildings, including calcium-carbonate marble, acidifies lakes and kills fish, and mobilizes soil aluminum that damages roots. Chemical remedies wash flue gas with limestone slurry to fix \chem{SO_2} by reaction with calcium carbonate, or burn low-sulfur coal. Elemental problems require elemental answers.

\section{Where This Account Fails}

Active, idle, fiery, and mild are useful shorthand but have boundaries; crossing them misleads.

First, activity and idleness depend on \textbf{conditions}, not permanent elemental personalities. Nitrogen watches at room temperature yet reacts in lightning, Haber's pressure vessel, and beside nitrogen-fixing enzymes. Oxygen makes iron burn fiercely in a pure-gas bottle but takes weeks to rust a nail in ordinary air. Possibility and rate are separate energy and kinetics accounts (Chapters 27 and 29). Chapter 13 already showed why inert cannot mean eternally unreactive.

Second, calling an element beneficial or harmful has no scientific meaning. Ozone protects overhead and injures lungs below; phosphorus sustains bones and ATP but can devastate lakes; trace hydrogen sulfide signals smell and bodily messages while high levels kill. Even excessive oxygen damages premature infants' retinas. \textbf{Location, dose, and chemical form} jointly determine role. Carry this principle into every wellness and detox advertisement.

Third, these accounts approximate typical conditions. Real systems contain competing reactions. Phosphorus does not always limit algae; some lakes and coastal waters are nitrogen-limited. Sulfuric-versus-nitric contributions to acid rain vary with regional energy use. Models prioritize which variable to inspect rather than supply universal answers.

\section{Look Again at the Match}

Replay the opening second through particle glasses.

Your hand creates friction, rapidly raising local temperature. A little red phosphorus on the striking strip converts into a more reactive form and ignites. Sparks reach the head, where heated potassium chlorate decomposes, releasing concentrated \chem{O_2}, a much faster-acting package than air's 21\%. It attacks sulfur:
\[\chem{S} + \chem{O_2} \rightarrow \chem{SO_2}\]
Sulfur's heat lights the stem, while the pungent smell is \chem{SO_2}. Flame travels along wood because its fibers, mainly cellulose containing carbon, hydrogen, and oxygen, keep reaching temperatures that pay the startup fee and react with oxygen.

Nitrogen, four fifths of the air, flows past throughout. It does not refuse participation; the triple-bond fee exceeds most of the match's budget. It simply heats, expands, and leaves, producing at most tiny amounts of nitrogen oxides in the hottest spots.

One match, four elements, each playing according to bond energy and conditions. Chemistry means seeing that card table clearly.

\section{Following the Chain Onward}

For the first time, elemental neighborhoods have left the table for air, bones, farmland, and chimneys. Three trails continue. Why does Haber's iron catalyst provide another route, and how do iron, vanadium, and manganese arrange molecules on surfaces? Chapter 15 visits that district. How do nitrogen and phosphorus cycle among rocks, oceans, and life, and how has humanity transformed those cycles? Chapter 16 draws the material map. How cells earn and spend ATP's small change awaits Volume II's Chapter 54 on respiration.

\begin{tiaozhan}
\textbf{Why Do Ammonia Plants Choose Conditions That Please Neither Extreme?} (Optional for the main story.) Ammonia synthesis releases heat and reduces gas molecules from four to two. Thermodynamics favors low temperature and high pressure. Why not room temperature and maximum pressure? Kinetics intervenes: cooling makes the \chem{N_2} startup fee harder to pay, slowing production below one tank a day. Pressure is constrained by steel strength and compression energy; doubling it more than doubles equipment cost and accident risk. Industry compromises around 200 atmospheres and $400\sim 500\,^{\circ}\mathrm{C}$ with iron catalysis, separating unreacted \chem{N_2} and \chem{H_2} for recycling. Each pass converts only part, but enough passes give high overall conversion. Equilibrium, rate, cost, and safety pull against one another. The optimum is an engineering compromise, not the maximum of one law, as real chemistry nearly always is.
\end{tiaozhan}

\section*{Try It Yourself}

\begin{sikaoti}
\item Draw shared electron pairs in \chem{O_2} and \chem{N_2}, beside bars for their splitting energies. Use this to explain why nitrogen, not oxygen, preserves chips.
\item Draw an ozone-location/role comparison with ground and heights above 20 kilometers on a horizontal axis. Mark ozone's work and human implications at each. Why is asking whether ozone is good or bad the wrong question?
\item Haber's pressure experiment overturned Nernst's pessimism. Which variable did Nernst's reasoning omit, and what does this teach about theory versus experiment?
\item Trace sulfur from coal through chimney, atmosphere, cloud water, and a lake's fish in a material-flow diagram. Label every arrow's physical movement or chemical reaction and its products.
\item A $189\ \mathrm{m^3}$ classroom contains 21\% oxygen. Forty-five students each consume $0.25\ \mathrm{L}$ per minute. Without ventilation, how long to consume 10\% of the oxygen? Why would real stuffiness arrive much sooner? (Recall carbon dioxide.)
\item An advertisement says an ozone generator disinfects bedrooms, adds oxygen, and improves health. Identify at least two problems and the evidence for your judgment.
\end{sikaoti}
